\section{扩散模型与流模型的等价关系}

本节是本讲的核心理论部分，讲者详细推导了扩散模型和流模型之间的精确数学等价关系。

\subsection{从变分扩散模型出发}

回忆变分扩散模型（Variational Diffusion Model, VDM）中，前向过程定义为：
$$q(\mathbf{x}_t \mid \mathbf{x}_0) = \mathcal{N}(\mathbf{x}_t; \alpha_t \mathbf{x}_0, \sigma_t^2 \mathbf{I})$$

在 Flow Matching 的语境下，将这个前向过程重新解读为条件流映射：
$$\psi_t(\mathbf{x}_0 \mid \mathbf{x}_1) = \sigma_t \cdot \mathbf{x}_0 + \alpha_t \cdot \mathbf{x}_1$$

\begin{importantbox}{核心洞察：扩散模型是流模型的特殊情况}
当条件流映射定义为上述线性高斯形式，且 $\alpha_t, \sigma_t$ 函数与变分扩散模型中的调度函数完全一致时，扩散模型可以被视为流模型的一种\textbf{特殊情况}。换言之，流模型是扩散模型的\textbf{推广}。两者的区别仅在于：
\begin{itemize}
    \item \textbf{扩散模型}：先定义概率路径（前向加噪过程），再推导出向量场和轨迹
    \item \textbf{流模型}：先定义条件流映射（轨迹），再推导出向量场和概率路径
\end{itemize}
\end{importantbox}

\subsection{Tweedie 公式的再现}

讲者在此处引入了一个重要的计算工具——Tweedie 公式。给定中间数据点 $\mathbf{x}_t = \sigma_t \mathbf{x}_0 + \alpha_t \mathbf{x}_1$，其中 $\mathbf{x}_0 \sim \mathcal{N}(\mathbf{0}, \mathbf{I})$, $\mathbf{x}_1 \sim p_{\text{data}}$，我们可以计算 $\mathbf{x}_1$ 的后验期望：

$$\mathbb{E}[\mathbf{x}_1 \mid \mathbf{x}_t] = \frac{\mathbf{x}_t - \sigma_t \cdot \mathbb{E}[\mathbf{x}_0 \mid \mathbf{x}_t]}{\alpha_t} = \frac{\mathbf{x}_t + \sigma_t^2 \nabla_{\mathbf{x}_t} \log p_t(\mathbf{x}_t)}{\alpha_t}$$

其中：
\begin{itemize}
    \item $\nabla_{\mathbf{x}_t} \log p_t(\mathbf{x}_t)$ 是分数函数（Score Function）
    \item 这与 Lecture 5 中介绍的 Tweedie 公式完全一致，只是使用了 Flow Matching 的符号约定
\end{itemize}

\begin{knowledgebox}{Tweedie 公式的直觉}
Tweedie 公式的核心思想是：给定一个被噪声污染的观测 $\mathbf{x}_t$，原始信号 $\mathbf{x}_1$ 的后验均值可以通过分数函数来计算。分数函数指向对数概率密度增长最快的方向，即``数据更密集''的方向。在高斯似然模型中，这等价于从当前位置向后验均值方向移动。
\end{knowledgebox}

\subsection{速度场与分数函数的关系}

这是本讲最重要的推导之一。从条件向量场出发：

$$u_t(\mathbf{x} \mid \mathbf{x}_1) = \dot{\sigma}_t \cdot \mathbf{x}_0 + \dot{\alpha}_t \cdot \mathbf{x}_1$$

将 $\mathbf{x}_0$ 用 $\mathbf{x}_t$ 和 $\mathbf{x}_1$ 表示：$\mathbf{x}_0 = \frac{\mathbf{x}_t - \alpha_t \mathbf{x}_1}{\sigma_t}$，代入得到：

$$u_t(\mathbf{x}_t \mid \mathbf{x}_1) = \frac{\dot{\sigma}_t}{\sigma_t}(\mathbf{x}_t - \alpha_t \mathbf{x}_1) + \dot{\alpha}_t \mathbf{x}_1$$

$$= \frac{\dot{\sigma}_t}{\sigma_t} \mathbf{x}_t + \left(\dot{\alpha}_t - \frac{\dot{\sigma}_t \alpha_t}{\sigma_t}\right) \mathbf{x}_1$$

对 $\mathbf{x}_1$ 取期望（边际化），并利用 Tweedie 公式，得到边际向量场：

\begin{importantbox}{速度场与分数函数的等价公式}
$$u_t(\mathbf{x}_t) = \frac{\dot{\sigma}_t}{\sigma_t} \mathbf{x}_t + \left(\dot{\alpha}_t - \frac{\dot{\sigma}_t \alpha_t}{\sigma_t}\right) \mathbb{E}[\mathbf{x}_1 \mid \mathbf{x}_t]$$

将 Tweedie 公式代入 $\mathbb{E}[\mathbf{x}_1 \mid \mathbf{x}_t]$，可以建立\textbf{速度场}与\textbf{分数函数}之间的等价关系：
$$u_t(\mathbf{x}_t) = \frac{\dot{\sigma}_t}{\sigma_t} \mathbf{x}_t + \left(\dot{\alpha}_t - \frac{\dot{\sigma}_t \alpha_t}{\sigma_t}\right) \cdot \frac{\mathbf{x}_t + \sigma_t^2 \nabla_{\mathbf{x}_t} \log p_t(\mathbf{x}_t)}{\alpha_t}$$

这意味着：给定任意中间数据点 $\mathbf{x}_t$ 和时间 $t$，只要知道分数函数就能计算速度场，反之亦然。
\end{importantbox}

\subsection{速度场与噪声预测的关系}

利用分数函数与噪声预测之间的已知关系 $\nabla_{\mathbf{x}_t} \log p_t(\mathbf{x}_t) = -\frac{\boldsymbol{\epsilon}}{\sigma_t}$，可以进一步建立速度场与噪声预测网络之间的等价关系：

$$u_t(\mathbf{x}_t) = \frac{\dot{\alpha}_t}{\alpha_t} \mathbf{x}_t - \left(\dot{\alpha}_t \frac{\sigma_t}{\alpha_t} - \dot{\sigma}_t\right) \boldsymbol{\epsilon}_\theta(\mathbf{x}_t, t)$$

其中 $\boldsymbol{\epsilon}_\theta$ 是扩散模型中训练好的噪声预测网络。

\begin{warningbox}{速度预测 vs.\ 噪声预测：不只是换个输出头}
虽然速度场和噪声场在数学上可以精确互转，但在\textbf{训练}时选择预测哪个量会影响梯度的数值行为。近年来的实践表明，直接训练速度预测（v-prediction）在某些调度下比噪声预测（$\epsilon$-prediction）具有更好的数值稳定性，特别是在 $t$ 接近端点时。这也是 SD3 和 Flux 选择 v-prediction 的原因之一。
\end{warningbox}

\subsection{与 Probability Flow ODE 的一致性}

讲者进一步指出，从 Flow Matching 推导出的向量场与从 SDE 框架推导出的 Probability Flow ODE（PF-ODE）是完全一致的。回忆 PF-ODE 的一般形式：

$$d\mathbf{x} = \left[f(\mathbf{x}, t) - \frac{1}{2}g(t)^2 \nabla_\mathbf{x} \log p_t(\mathbf{x})\right] dt$$

其中 $f(\mathbf{x}, t)$ 和 $g(t)$ 分别是 SDE 的漂移项和扩散系数，它们与 $\alpha_t, \sigma_t$ 的关系为：
$$f(\mathbf{x}, t) = \frac{\dot{\alpha}_t}{\alpha_t} \mathbf{x}, \quad g(t)^2 = \frac{d\sigma_t^2}{dt} - 2\frac{\dot{\alpha}_t}{\alpha_t}\sigma_t^2$$

将这些代入 PF-ODE，可以验证其结果与 Flow Matching 推导的向量场\textbf{完全相同}。

\begin{importantbox}{三种等价视角的统一}
到此为止，我们建立了以下三种视角的完全等价关系（在线性高斯条件路径下）：
\begin{enumerate}
    \item \textbf{扩散模型视角}：学习分数函数 $\nabla_\mathbf{x} \log p_t(\mathbf{x})$ 或噪声 $\boldsymbol{\epsilon}$
    \item \textbf{Flow Matching 视角}：学习向量场/速度 $u_t(\mathbf{x})$
    \item \textbf{PF-ODE 视角}：通过 SDE 到 ODE 的转换获得确定性采样路径
\end{enumerate}
三者在数学上可精确互转，区别在于训练时预测哪个量、以及采样时使用哪种求解策略。
\end{importantbox}

\subsection{本章小结}

本节建立了扩散模型与流模型之间的精确数学等价关系。核心结论是：在线性高斯条件流映射下，扩散模型是流模型的特殊情况。速度场、分数函数、噪声预测三者可以自由互转。这一等价关系不仅具有理论意义，更在实际中意味着：一个预训练好的扩散模型可以直接被``重新解读''为流模型来使用。

